\documentclass[11pt]{article}

\usepackage[a4paper,margin=1in]{geometry}
\usepackage{amsmath,amssymb,bm}
\usepackage{booktabs}
\usepackage{graphicx}
\usepackage{hyperref}
\usepackage{enumitem}

\hypersetup{
  colorlinks=true,
  linkcolor=blue,
  urlcolor=blue,
  citecolor=blue
}

\title{Calculating Propensities from a Core-Hole Potential}


\begin{document}

\section{Core-Hole Potential}

Instead of an explicit calculating of the overlap term, $S_{pq}$ - which depends on the final relaxed orbital $\phi_p$ and initial orbital $\psi_q$, these can also be obtained in terms of the core-hole potential acting on the initial wave function. In first order perturbation theory the overlap is then,

\begin{equation}
  S_{pq} = \langle\phi_p|\psi_q \rangle \approx
      \frac{\langle\psi_p|\hat{v}_{\mathrm{ch}}|\psi_q\rangle}{\epsilon_p - \epsilon_q}
\label{eq:final_propensity_rule}
\end{equation}

with $\psi_p$ and $\psi_q$ the orbitals involved in shake-down where $p \neq q$, and $\epsilon$ defines their MO energies. We wish to compute the matrix-element,
\begin{equation}
    \langle \psi_p|\hat{v}_{\mathrm{ch}}|\psi_q \rangle \approx \int  dr \hspace{2mm} \psi_p^*(\mathbf{r}) v_{ch}(\mathbf{r})\psi_q(\mathbf{r}),
\end{equation}
with,
\begin{equation}
    v_{ch}(\mathbf{r}) = - \int d\mathbf{r}^{\prime} \hspace{2mm} \frac{\rho_c (\mathbf{r}^{\prime})}{|\mathbf{r}-\mathbf{r}^{\prime}|},
\end{equation}    
where $\rho_c$ is the density of the core orbital hole. We can expand each molecular orbital in the AO basis,
\begin{equation}
  \psi_p(\mathbf{r}) = \sum_{\alpha} C_{\alpha p} \chi_\alpha(\mathbf{r}),
\end{equation}
giving the AO representation of the matrix element
\begin{equation}
  \langle \psi_p|\hat{v}_{\mathrm{ch}}|\psi_q\rangle
  \approx
  \sum_{\alpha,\beta}
  C_{\alpha p}^{*}
  V_{\alpha\beta}
  C_{\beta q},
  \label{eq:local_core_hole_ao}
\end{equation}
where $V_{\alpha\beta}$ is the AO representation of the core-hole potential $v_{ch}$,

\begin{equation}
  V_{\alpha\beta}
  =
  -\iint d\mathbf r\,d\mathbf r'
  \frac{
  \chi_\alpha^{*}(\mathbf r)
  \chi_\beta(\mathbf r)
  \rho_c(\mathbf r')
  }
  {|\mathbf r-\mathbf r'|}.
  \label{eq:core_hole_ao_integral}
\end{equation}
Expanding the core-orbital $\psi_c$ (in the density $\rho_c$) in an AO basis contribution,
\begin{equation}
  \psi_{c}(\mathbf r)
  =
  \sum_{\gamma}
  C_{\gamma c}\chi_\gamma(\mathbf r),
\end{equation}
we arrive at an expression for the core-hole potential as a regular two-electron integral, 
\begin{equation}
  V_{\alpha\beta}
  =
  -
  \sum_{\gamma\delta}
  C_{\gamma c}^{*}
  C_{\delta c}
  (\alpha\beta|\gamma\delta),
  \label{eq:core_hole_eri}
\end{equation}
where,

\begin{equation}
  (\alpha\beta|\gamma\delta)
  =
  \iint d\mathbf r\,d\mathbf r'
  \frac{
  \chi_\alpha^*(\mathbf r)
  \chi_\beta(\mathbf r)
  \chi_\gamma^*(\mathbf r')
  \chi_\delta(\mathbf r')
  }
  {|\mathbf r-\mathbf r'|}.
\end{equation}
Thus giving the total matrix element in the AO representation,
\begin{equation}
  \boxed{
  \langle \psi_p|\hat{v}_{\mathrm{ch}}|\psi_q\rangle
  \approx
  -
  \sum_{\alpha,\beta}
  \sum_{\gamma,\delta}
  C_{\alpha p}^{*}
  C_{\beta q}
  C_{\gamma c}^{*}
  C_{\delta c}
  (\alpha\beta|\gamma\delta)
  }
  \label{eq:core_hole_final}
\end{equation}
%
In MO representation this is simply the four-center MO integral
\begin{equation}
    \boxed{
  \langle \psi_p|\hat{v}_{\mathrm{ch}}|\psi_q\rangle
  \approx
  -(p q | c c)
  }
\end{equation}

Thus,
\begin{equation}
    S_{pq} \approx \frac{(pq|cc)}{\epsilon_q - \epsilon_p}
\end{equation}

\end{document}